% !TEX program = xelatex
\documentclass[10pt,a4paper]{article}
\usepackage[a4paper,margin=0pt]{geometry}
\usepackage{fontspec}
\usepackage{polyglossia}
\usepackage[dvipsnames]{xcolor}
\usepackage{tikz}
\usepackage{eso-pic}
\usepackage{enumitem}
\usepackage{tabularx}
\usepackage{array}
\usepackage{calc}
\usepackage{graphicx}
\usepackage{textcomp}
\usepackage[hidelinks]{hyperref}
\usepackage[final]{microtype}
\usepackage{ragged2e}

\defaultfontfeatures{Ligatures=TeX}
\setmainfont{Rubik-Regular.ttf}[
  Path = fonts/,
  BoldFont = Rubik-Bold.ttf,
  ItalicFont = Rubik-Italic.ttf,
  BoldItalicFont = Rubik-BoldItalic.ttf
]
\newfontfamily\RubikRegular{Rubik-Regular.ttf}[
  Path = fonts/,
  BoldFont = Rubik-Medium.ttf,
  ItalicFont = Rubik-Italic.ttf,
  BoldItalicFont = Rubik-BoldItalic.ttf
]
\newfontfamily\RubikLight{Rubik-Light.ttf}[
  Path = fonts/,
  ItalicFont = Rubik-Italic.ttf
]
\newfontfamily\TemplateFont{cmunrm.otf}[
  Path = fonts/,
  BoldFont = cmunbx.otf,
  ItalicFont = cmunti.otf,
  BoldItalicFont = cmunbi.otf
]
\newfontfamily\cyrillicfont{Rubik-Regular.ttf}[
  Script = Cyrillic,
  Path = fonts/,
  BoldFont = Rubik-Bold.ttf,
  ItalicFont = Rubik-Italic.ttf,
  BoldItalicFont = Rubik-BoldItalic.ttf
]

\pagestyle{empty}
\setlength{\parindent}{0pt}
\setlength{\parskip}{0pt}
\setlength{\tabcolsep}{0pt}
\setlength{\fboxsep}{0pt}
\setlength{\emergencystretch}{2em}
\raggedbottom
\urlstyle{same}
\hypersetup{colorlinks=true, urlcolor=black, linkcolor=black}

\definecolor{sidebarbg}{HTML}{244869}
\definecolor{accentblue}{HTML}{0B77F3}
\definecolor{textgray}{HTML}{555555}
\definecolor{rulegray}{HTML}{B7B7B7}
\definecolor{sidebarline}{HTML}{DCE6F0}
\definecolor{softwhite}{HTML}{F7FAFC}

\newlength{\sidebarwidth}
\setlength{\sidebarwidth}{0.355\paperwidth}
\newlength{\mainwidth}
\setlength{\mainwidth}{\paperwidth-\sidebarwidth}
\newlength{\sidebarinnerwidth}
\setlength{\sidebarinnerwidth}{\sidebarwidth-2.12cm}
\newlength{\maininnerwidth}
\setlength{\maininnerwidth}{\mainwidth-1.72cm}
\newlength{\pagecontentheight}
\setlength{\pagecontentheight}{\paperheight-1.05cm}
\newlength{\sidebarinnerheight}
\setlength{\sidebarinnerheight}{\pagecontentheight-0.84cm}
\newlength{\maininnerheight}
\setlength{\maininnerheight}{\pagecontentheight-0.74cm}
\newlength{\maindatewidth}
\setlength{\maindatewidth}{2.65cm}
\newlength{\maingap}
\setlength{\maingap}{0.30cm}

\newlength{\avatarheight}
\setlength{\avatarheight}{5.33cm}

\newcommand{\atsbackground}{%
  \AddToShipoutPictureBG*{%
    \AtPageLowerLeft{\color{sidebarbg}\rule{\sidebarwidth}{\paperheight}}%
  }%
}

\newcommand{\TemplateText}[1]{{\TemplateFont #1}}
\newcommand{\MetaLabel}[1]{{\TemplateFont\fontsize{8.55}{8.9}\selectfont #1}}

\newcommand{\SidebarAvatar}{%
  {\centering\includegraphics[height=\avatarheight,keepaspectratio]{data/avatar.png}\par}
  \vspace{0.30cm}%
}
\newcommand{\CVName}[1]{%
  {\centering\TemplateFont\fontsize{25.0}{24.4}\selectfont\bfseries\MakeUppercase{#1}\par}
  \vspace{0.66cm}
}

\newcommand{\SidebarSection}[1]{%
  {\TemplateFont\fontsize{12.2}{12.7}\selectfont\MakeUppercase{#1}\par}
  \vspace{-0.30cm}
  {\color{sidebarline}\rule{\linewidth}{0.48pt}\par}
  \vspace{0.10cm}
}

\newcommand{\SidebarRoleTitle}[1]{{\RubikRegular\fontsize{9.9}{10.55}\selectfont\bfseries #1\par}}
\newcommand{\SidebarOrg}[1]{{\RubikRegular\fontsize{9.15}{9.8}\selectfont #1\par}}
\newcommand{\SidebarDate}[1]{{\RubikRegular\fontsize{8.88}{9.4}\selectfont #1\par}}
\newcommand{\SidebarBody}[1]{{\RubikLight\fontsize{8.95}{9.7}\selectfont #1\par}}

\newenvironment{sidebarbullets}{%
  \begin{itemize}[leftmargin=1.08em,topsep=0.08em,itemsep=0.05em,parsep=0pt,partopsep=0pt]
    \RubikLight\fontsize{8.72}{9.55}\selectfont\color{softwhite}
}{\end{itemize}}

\newcommand{\SidebarEntrySep}{\vspace{0.30cm}}
\newcommand{\SidebarSectionSep}{\vspace{0.68cm}}

\newcommand{\EducationEntry}[4]{%
  {\fontsize{10.15}{10.65}\selectfont\bfseries #1\par}
  \vspace{0.04cm}
  {\fontsize{9.1}{9.65}\selectfont #2\par}
  \vspace{0.05cm}
  {\fontsize{8.9}{9.4}\selectfont #3\ \textbar\ #4\par}
}

\newcommand{\LanguageEntry}[2]{%
  \begin{tabular*}{\linewidth}{@{}>{\fontsize{9.7}{10.15}\selectfont}l@{\extracolsep{\fill}}>{\fontsize{9.1}{9.55}\selectfont}r@{}}
    #1 & #2
  \end{tabular*}\par
}

\newcommand{\SkillLine}[1]{{\fontsize{9.2}{9.75}\selectfont #1\par}}

\newcommand{\InterestTitle}[1]{{\RubikRegular\fontsize{10.35}{10.95}\selectfont\bfseries #1\par\vspace{0.03cm}}}
\newcommand{\InterestBody}[1]{{\RubikLight\fontsize{8.75}{9.6}\selectfont #1\par}}

\newcommand{\TopRole}[1]{{\fontsize{13.8}{14.15}\selectfont\color{accentblue}#1\par}}
\newcommand{\ContactLine}[1]{{\fontsize{8.25}{8.7}\selectfont\color{textgray}#1\par}}
\newcommand{\HeaderGap}{\vspace{0.46cm}}

\newcommand{\MainSection}[1]{%
  {\TemplateFont\fontsize{13.6}{14.1}\selectfont\MakeUppercase{#1}\par}
  \vspace{-0.30cm}
  {\color{rulegray}\rule{\linewidth}{0.46pt}\par}
  \vspace{0.12cm}
}

\newcommand{\SummaryText}[1]{
  {\fontsize{9.9}{10.5}\selectfont\color{black!85}#1\par}
  \vspace{0.35cm}
}

\newcommand{\MainEntry}[5]{%
  \begin{tabular*}{\linewidth}{@{}p{\dimexpr\linewidth-\maindatewidth-\maingap\relax}@{\hspace{\maingap}}>{\raggedleft\arraybackslash}p{\maindatewidth}@{}}
    {\fontsize{10.75}{11.1}\selectfont #1} & {\fontsize{9.2}{9.6}\selectfont #2}
  \end{tabular*}\par
  \vspace{-0.30cm}
  \begin{tabular*}{\linewidth}{@{}p{\dimexpr\linewidth-\maindatewidth-\maingap\relax}@{\hspace{\maingap}}>{\raggedleft\arraybackslash}p{\maindatewidth}@{}}
    {\fontsize{9.9}{10.3}\selectfont\color{accentblue}#3} & {\fontsize{8.95}{9.35}\selectfont\raggedleft\color{textgray}#4}
  \end{tabular*}\par
  \vspace{0.02cm}
  #5
  \vspace{0.35cm}
}

\newenvironment{mainbullets}{%
  \begin{itemize}[leftmargin=1.0em,topsep=0.02em,itemsep=0.04em,parsep=0pt,partopsep=0pt]
    \fontsize{9.3}{9.8}\selectfont\color{black!82}
}{\end{itemize}}

\newcommand{\CoursesSection}[1]{%
  {\TemplateFont\fontsize{13.6}{14.1}\selectfont\MakeUppercase{#1}\par}
  \vspace{-0.30cm}
  {\color{rulegray}\rule{\linewidth}{0.46pt}\par}
  \vspace{0.12cm}
}

\newcommand{\CourseItem}[2]{%
  {\fontsize{9.9}{10.3}\selectfont\color{accentblue}#1\par}
  \vspace{0.02cm}
  {\fontsize{9.15}{9.6}\selectfont\color{black!82}#2\par}
  \vspace{0.30cm}
}

\newcommand{\RightFooter}[1]{\vfill{\hfill\fontsize{8.0}{8.0}\selectfont\color{textgray}#1\par}}
\newcommand{\LeftFooter}[1]{\vfill{\fontsize{8.0}{8.0}\selectfont\color{sidebarline}#1\par}}

\setmainlanguage{english}
\setotherlanguage{russian}

\begin{document}
\atsbackground
\noindent\makebox[\paperwidth][l]{%
\begin{minipage}[t][\pagecontentheight]{\sidebarwidth}
  \vspace*{0.46cm}
  \hspace*{1.06cm}\begin{minipage}[t][\sidebarinnerheight]{\sidebarinnerwidth}
    \color{softwhite}\RaggedRight
    \SidebarAvatar
    \CVName{Robert\\Bitlev}

    \SidebarSection{Achievements}
    \SidebarRoleTitle{Science Olympiads\\and competitions}
    \vspace{0.05cm}
    \begin{sidebarbullets}
      \item Gold medalist at the International Junior Science Olympiad (IJSO) in 2023 and 2024.
      \item Winner and prize-winner of the final stage of the All-Russian Olympiad for School Students in Physics.
      \item Winner of the Moscow Olympiad in Physics in 2023 and 2024.
      \item Final-stage participant in the IChem Prize in 2024.
    \end{sidebarbullets}

    \SidebarSectionSep
    \SidebarSection{Education}
    \EducationEntry{B.Sc. in Computer Science}{St. Petersburg State University}{2025--2029}{St. Petersburg}
    \SidebarEntrySep
    \EducationEntry{Secondary general education}{Lyceum No.\,533}{2016--2025}{GPA: 5.00/5.00}

    \SidebarSectionSep
    \SidebarSection{Languages}
    \LanguageEntry{English}{C1}
    \SidebarEntrySep
    \LanguageEntry{German}{B2--C1}
    \SidebarEntrySep
    \LanguageEntry{Russian}{Native}

    \SidebarSectionSep
    \SidebarSection{Skills}
    \SkillLine{Physics \textperiodcentered\ Mathematics}
    \SidebarEntrySep
    \SkillLine{Python \textperiodcentered\ Data analysis}
    \SidebarEntrySep
    \SkillLine{Arduino \textperiodcentered\ Research projects}
    \SidebarEntrySep
    \SkillLine{Radio electronics}

    \SidebarSectionSep
    \SidebarSection{Interests}
    \InterestTitle{Natural sciences\\and technology}
    \InterestBody{Interested in interdisciplinary problems at the intersection of physics, mathematics, engineering, and software tools.}
    \LeftFooter{}
  \end{minipage}
\end{minipage}%
\begin{minipage}[t][\pagecontentheight]{\mainwidth}
  \vspace*{0.70cm}
  \hspace*{0.66cm}\begin{minipage}[t][\maininnerheight]{\maininnerwidth}
    \TopRole{AI360 Student \textbar\ Physics and Electronics}
    \vspace{0.09cm}
    \ContactLine{{\MetaLabel{Tel.:}}\ \href{tel:+79215747178}{+7 (921) 574-71-78} \textbar\ {\MetaLabel{Email:}}\ \href{mailto:robertbitlev09@gmail.com}{robertbitlev09@gmail.com}}
    \vspace{0.03cm}
    \ContactLine{{\MetaLabel{GitHub:}}\ \href{https://github.com/robrob09/}{robrob09} \textbar\ {\MetaLabel{LinkedIn:}}\ \href{https://www.linkedin.com/in/robert-bitlev-329662377/}{robert-bitlev}}
    \HeaderGap

    \MainSection{Summary}
    \SummaryText{AI360 student at the Faculty of Mathematics and Computer Science with a strong background in physics, mathematics, and electronics. Completed olympiad, research, and interdisciplinary projects that combine experimental work, hardware prototyping, data analysis, and software tools.}

    \vspace{0.24cm}
    \MainSection{Projects}
    \MainEntry
      {Radio: Bridge over Waves}
      {Spring 2024}
      {Speech capture and transcription system}
      {}
      {\begin{mainbullets}
        \item Built 2 Arduino-based modules: a transmitter with a microphone, encryption, and a 2.4 GHz radio link, and a receiver that streams data to a server over USB.
        \item Implemented a 3-stage Python pipeline to reconstruct missing samples, save WAV audio, and convert speech to text via a cloud API and an offline recognizer.
      \end{mainbullets}}

    \MainEntry
      {Silver nanoparticle thin films}
      {2023--2024}
      {SERS films for sensors}
      {}
      {\begin{mainbullets}
        \item Synthesized silver nanoparticles by hydroxylamine reduction, calculated 2,133 uL of colloid for a 7.6 cm2 TTF interface, and transferred SERS films onto about 0.8 x 0.9 cm silanized substrates.
        \item Characterized samples with UV-Vis and AFM in 1 x 1 um scans, estimated mean particle diameter at about 79 nm, and observed film defects no larger than 300 nm.
      \end{mainbullets}}

    \MainEntry
      {Brain Train}
      {2023--2024}
      {Interactive memory trainer on Arduino Nano}
      {}
      {\begin{mainbullets}
        \item Designed an Arduino Nano device with random LED-sequence generation, OLED feedback, hint logic, and button input for attention and short-term memory training.
        \item Integrated 4 Arduino C++ modules - game logic, button handling, result display, and timing control - to make the training loop more stable and responsive.
      \end{mainbullets}}

    \MainEntry
      {Modern Systematics of Protists}
      {2021--2022}
      {Research and taxonomy catalog}
      {}
      {\begin{mainbullets}
        \item Built a SQLite reference catalog of 73 organisms across 6 tables, reflecting recent taxonomy revisions and supporting filtering by 4 biological traits plus Russian and Latin name search.
        \item Added database updates through text and CSV files with 3 commands (ADD, CHANGE, DELETE) and an in-app glossary of 33 biological terms for continued study and maintenance.
      \end{mainbullets}}

    \CoursesSection{Additional Education / Courses}
    \CourseItem{Physics Center}{Presidential Physics and Mathematics Lyceum No.\,239, 2019--2025}
    \CourseItem{School of practical programming and data analysis}{HSE University, 2024--2025}
    \CourseItem{Creative radio electronics}{Academy of Digital Technologies, 2023--2024}
    \CourseItem{Programming courses}{Competitive programming, 2023--2025}
    \CourseItem{Programming school}{Fundamentals of Industrial Programming, 2021--2022}
    \CourseItem{Experimental physics and non-standard problem solving}{Physics Center, 2019--2025}
    \RightFooter{}
  \end{minipage}
\end{minipage}}
\end{document}
